% Labeling-algorithm figure: how one stability label is produced.
%
% Panel (a) is a schematic of the procedure in Sec. 3.
% Panels (b) and (c) are real data:
%   (b) mean-over-branches log divergence, regenerated for
%       rc_PickPlaceCounterToCabinet demo 20 by replaying the labeler's exact
%       call sequence and RNG stream (seed = demo index = 20) and keeping the
%       per-step curves that the production labeler discards.  The recomputed
%       slopes reproduce the published labels to <1e-4, so the lambda values
%       quoted here are the same numbers plotted in Fig. 1.
%   (c) histogram of lambda_task over all 23,724 stamps of that task, with the
%       automatic threshold delta = .0649 and the 5.0% of stamps above it.
% Regeneration script: scratchpad/dump_divergence.py

\begin{figure}[t]
\centering

%% ---------------------------------------------------------------- panel (a)
\begin{tikzpicture}[font=\footnotesize,
  stp/.style={circle, draw=cline, fill=white, inner sep=0.5pt,
              minimum size=2.7mm, font=\tiny, line width=0.4pt},
  sbox/.style={draw=cline, fill=cbox, rounded corners=1.6pt, align=center,
               inner sep=2.5pt, font=\tiny, line width=0.5pt},
  ar/.style={-latex, cline, line width=0.5pt},
]
% tau = 0 at x=2.4, tau = 24 at x=9.2 (0.2833 cm per step); nominal at y=1.75
\node[sbox, text width=10mm] (s0) at (1.05,1.75) {restore\\$s_{t_0}$};
\draw[ar] (s0.east) -- (2.36,1.75);
\node[stp] at (0.22,1.75) {1};

% the N=8 perturbed branches: same recorded actions, different first action
\draw[cline!60, line width=0.5pt] (2.4,1.75) .. controls (5.0,1.918) and (7.4,2.254) .. (9.2,2.87);
\draw[cline!60, line width=0.5pt] (2.4,1.75) .. controls (5.0,1.879) and (7.4,2.137) .. (9.2,2.61);
\draw[cline!60, line width=0.5pt] (2.4,1.75) .. controls (5.0,1.839) and (7.4,2.016) .. (9.2,2.34);
\draw[cline!60, line width=0.5pt] (2.4,1.75) .. controls (5.0,1.798) and (7.4,1.894) .. (9.2,2.07);
\draw[cline!60, line width=0.5pt] (2.4,1.75) .. controls (5.0,1.711) and (7.4,1.633) .. (9.2,1.49);
\draw[cline!60, line width=0.5pt] (2.4,1.75) .. controls (5.0,1.671) and (7.4,1.511) .. (9.2,1.22);
\draw[cline!60, line width=0.5pt] (2.4,1.75) .. controls (5.0,1.628) and (7.4,1.386) .. (9.2,0.94);
\draw[cline!60, line width=0.5pt] (2.4,1.75) .. controls (5.0,1.578) and (7.4,1.233) .. (9.2,0.60);

% the nominal branch
\draw[cline, line width=0.9pt] (2.4,1.75) -- (9.2,1.75);
\fill[cline] (2.4,1.75) circle (1.3pt);

% what makes the branches differ
\node[font=\tiny, cunstable, anchor=south west, align=left, inner sep=1pt]
  (pert) at (2.95,2.45)
  {perturb $a_{t_0}$ only: $+\,\sigma_u\varepsilon$\\
   on the translation dims};
\draw[-latex, cunstable, line width=0.5pt] (3.05,2.43) -- (2.63,1.83);
\node[font=\tiny, cline, anchor=south west, inner sep=1pt]
  at (2.60,0.62) {identical actions in every branch after the first};

% the quantity that is measured
\draw[latex-latex, cline, line width=0.6pt] (7.5,1.75) -- (7.5,2.40);
\node[font=\tiny, fill=white, inner sep=1pt] at (7.5,2.08) {$d_n(\tau)$};
\node[stp] at (7.5,2.62) {4};

\node[stp] at (9.38,1.75) {2};
\node[font=\tiny, anchor=west, inner sep=1.5pt] at (9.52,1.75) {nominal};
\node[stp] at (9.38,2.87) {3};
\node[font=\tiny, anchor=west, inner sep=1.5pt, align=left]
  at (9.52,2.87) {$N{=}8$\\branches};

% tau axis
\draw[-latex, cline, line width=0.5pt] (2.4,0.30) -- (9.55,0.30);
\draw[cline, line width=0.5pt] (2.4,0.30) -- (2.4,0.21);
\draw[cline, line width=0.5pt] (4.1,0.30) -- (4.1,0.21);
\draw[cline, line width=0.5pt] (5.8,0.30) -- (5.8,0.21);
\draw[cline, line width=0.5pt] (7.5,0.30) -- (7.5,0.21);
\draw[cline, line width=0.5pt] (9.2,0.30) -- (9.2,0.21);
\node[font=\tiny, anchor=north, inner sep=2pt] at (2.4,0.21) {0};
\node[font=\tiny, anchor=north, inner sep=2pt] at (4.1,0.21) {6};
\node[font=\tiny, anchor=north, inner sep=2pt] at (5.8,0.21) {12};
\node[font=\tiny, anchor=north, inner sep=2pt] at (7.5,0.21) {18};
\node[font=\tiny, anchor=north, inner sep=2pt] at (9.2,0.21) {24};
\node[font=\tiny, anchor=west, inner sep=2pt] at (9.55,0.22) {$\tau$};

\node[anchor=north west, font=\footnotesize\bfseries] at (0.0,3.10) {(a)};
% the box is the full text width, and the (b)+(c) row below is stretched to
% the same width with \hfill, so the two rows share a left edge
\useasboundingbox (0.0,-0.08) rectangle (\textwidth,3.10);
\end{tikzpicture}

\vspace{2pt}

%% ---------------------------------------------------------------- panel (b)
\begin{tikzpicture}[font=\footnotesize]
\begin{axis}[
  scale only axis, width=4.5cm, height=2.6cm,
  xmin=1, xmax=24, ymin=-7.2, ymax=-3.0,
  xtick={1,6,12,18,24}, ytick={-7,-6,-5,-4},
  ticklabel style={font=\scriptsize},
  xlabel={$\tau$}, ylabel={mean $\log d_n(\tau)$},
  xlabel style={font=\scriptsize, yshift=3pt},
  ylabel style={font=\scriptsize, yshift=-4pt},
  axis line style={cline}, tick style={cline}, tick align=outside,
  every axis plot/.append style={line width=0.9pt},
  clip mode=individual,
]
  % t0=48: unstable, d grows 0.0018 m -> 0.0256 m over the window
  \addplot[cunstable] coordinates {
    (1,-6.3076) (2,-5.9927) (3,-6.2816) (4,-5.5934) (5,-5.3721) (6,-5.2209)
    (7,-5.1520) (8,-5.1344) (9,-5.1184) (10,-5.0567) (11,-5.1292) (12,-4.8690)
    (13,-4.5927) (14,-4.3482) (15,-4.1399) (16,-3.9602) (17,-3.8317)
    (18,-3.7741) (19,-3.7472) (20,-3.7123) (21,-3.6822) (22,-3.6773)
    (23,-3.6709) (24,-3.6637)
  };
  \addplot[cunstable, dashed, line width=0.6pt]
    coordinates {(1,-6.0568) (24,-3.2789)};
  % t0=10: stable, d stays at 0.0013-0.0018 m
  \addplot[cstable] coordinates {
    (1,-6.6133) (2,-6.4940) (3,-6.4810) (4,-6.4753) (5,-6.3353) (6,-6.2190)
    (7,-6.1582) (8,-6.1909) (9,-6.2181) (10,-6.2024) (11,-6.1560) (12,-6.1454)
    (13,-6.0897) (14,-6.0891) (15,-6.1115) (16,-6.1230) (17,-6.1617)
    (18,-6.2319) (19,-6.2902) (20,-6.3413) (21,-6.3524) (22,-6.3475)
    (23,-6.3614) (24,-6.3434)
  };
  \addplot[cstable, dashed, line width=0.6pt]
    coordinates {(1,-6.3393) (24,-6.2050)};
  \node[anchor=north west, font=\tiny, cunstable, align=left, inner sep=1pt]
    at (axis cs:1.4,-3.02) {unstable state, $\lambda{=}.121$};
  \node[anchor=south west, font=\tiny, cstable, align=left, inner sep=1pt]
    at (axis cs:1.4,-7.18) {stable state, $\lambda{=}.006$};
\end{axis}
\node[anchor=north west, font=\footnotesize\bfseries] at (-1.15,2.95) {(b)};
\end{tikzpicture}
\hfill
%% ---------------------------------------------------------------- panel (c)
\begin{tikzpicture}[font=\footnotesize]
\begin{axis}[
  scale only axis, width=4.5cm, height=2.6cm,
  xmin=-0.10, xmax=0.16, ymin=0, ymax=19.6,
  xtick={-0.1,0,0.1}, ytick={0,5,10,15},
  ticklabel style={font=\scriptsize},
  xlabel={$\lambda$ over all stamps of the task},
  ylabel={\% of stamps},
  xlabel style={font=\scriptsize, yshift=3pt},
  ylabel style={font=\scriptsize, yshift=-4pt},
  axis line style={cline}, tick style={cline}, tick align=outside,
  clip mode=individual,
]
  \fill[cband] (axis cs:0.0649,0) rectangle (axis cs:0.16,19.6);
  \addplot[ybar interval, draw=cstable, fill=cstable!25, line width=0.3pt]
    coordinates {
    (-0.10,0.510) (-0.09,0.691) (-0.08,0.784) (-0.07,0.906) (-0.06,1.290)
    (-0.05,1.960) (-0.04,3.267) (-0.03,6.171) (-0.02,10.538) (-0.01,16.051)
    (0.00,18.230) (0.01,13.674) (0.02,7.478) (0.03,5.105) (0.04,3.971)
    (0.05,3.216) (0.06,2.234) (0.07,1.467) (0.08,1.113) (0.09,0.577)
    (0.10,0.375) (0.11,0.181) (0.12,0.122) (0.13,0.046) (0.14,0.025)
    (0.15,0.013) (0.16,0)
  };
  \draw[cunstable, dashed, line width=0.7pt]
    (axis cs:0.0649,0) -- (axis cs:0.0649,19.6);
  \node[anchor=north east, font=\tiny, cunstable, inner sep=1pt]
    at (axis cs:0.0625,19.4) {$\delta$};
  \node[anchor=north east, font=\tiny, cunstable, align=right, inner sep=1pt]
    at (axis cs:0.158,14.0) {unstable:\\5.0\% of\\stamps};
\end{axis}
\node[anchor=north west, font=\footnotesize\bfseries] at (-1.15,2.95) {(c)};
\end{tikzpicture}

\caption{\textbf{How one stability label is produced.}
\textbf{(a)} The procedure at a single state, and the whole of it: restore the
simulator to $s_{t_0}$ (1); replay the $K{=}24$ recorded actions unperturbed to
get the nominal branch (2); repeat $N{=}8$ times with Gaussian noise added to
the translation components of the first action only, every branch replaying the
same recorded actions thereafter (3); and record the task-space distance
$d_n(\tau)$ between each branch and the nominal at every step (4). The label
$\lambda$ is the least-squares slope of the mean log distance against $\tau$.
Nothing is trained, no policy internals are read, and the policy is not involved
at all: the label describes the plant and the task geometry under the recorded
behavior.
\textbf{(b)} The fit, on two real states drawn from the same episode as
Figure~\ref{fig:framework}, with the least-squares line dashed. At the unstable
state the branches separate from 1.8 to 25.6 millimetres in the task metric
over the window, a factor of 14; at the stable state they stay between 1.3 and
1.8 millimetres,
and the fitted slope is twenty times smaller. The label is the slope,
not the final distance, so it does not reward states that merely start far
apart. Neither curve is a clean straight line, and the fit is a summary of the
window rather than a claim that growth is exactly exponential within it.
\textbf{(c)} Where the threshold comes from. Distribution of $\lambda$ over all
23{,}724 stamps of PickPlaceCounterToCabinet. The bulk sits at a task-specific
noise floor set by contact chatter and metric scale, so a fixed cutoff would not
transfer between tasks; $\delta$ is instead read off the below-median half of
this distribution as $P_{95}(\lvert\lambda\rvert)$, which puts 5.0\% of stamps
in the unstable class here. Stamps with $\lvert\lambda\rvert \le \delta$ are the
deadband and are excluded from classification training.}
\label{fig:labeling}
\end{figure}
